\documentclass[11pt, a4paper]{article}

\usepackage[T1]{fontenc}
\usepackage[utf8]{inputenc}
\usepackage{lmodern}
\usepackage[margin=22mm]{geometry}
\usepackage{amsmath, amssymb}
\usepackage{graphicx}
\usepackage{booktabs}
\usepackage{longtable}
\usepackage{array}
\usepackage[table]{xcolor}
\usepackage{tcolorbox}
\usepackage[ruled, vlined, linesnumbered]{algorithm2e}
\usepackage{listings}
\usepackage{caption}
\usepackage{enumitem}
\usepackage[hidelinks]{hyperref}
\usepackage{pdflscape}

\graphicspath{{figures/}}

\definecolor{codebg}{RGB}{247,247,247}
\definecolor{accent}{RGB}{43,108,176}
\definecolor{warn}{RGB}{192,72,40}

\lstset{
  basicstyle=\ttfamily\footnotesize,
  backgroundcolor=\color{codebg},
  frame=single, rulecolor=\color{black!20},
  breaklines=true, columns=fullflexible,
  keywordstyle=\color{accent}\bfseries,
  commentstyle=\color{black!55}\itshape,
  language=Matlab, showstringspaces=false,
  xleftmargin=2mm, framexleftmargin=2mm,
}

\newtcolorbox{keybox}[1]{colback=blue!4, colframe=accent, fonttitle=\bfseries,
  title=#1, boxrule=0.6pt, arc=2pt, left=3mm, right=3mm, top=2mm, bottom=2mm}
\newtcolorbox{warnbox}[1]{colback=red!3, colframe=warn, fonttitle=\bfseries,
  title=#1, boxrule=0.6pt, arc=2pt, left=3mm, right=3mm, top=2mm, bottom=2mm}

\title{\vspace{-14mm}\textbf{Aircraft Sizing}\\[2mm]
       \large The closure loop, and what a second level of fidelity buys}
\author{AOE 4065 --- Air Vehicle Design}
\date{}

% Long \texttt{} identifiers and file paths cannot hyphenate, so let TeX
% stretch the interword space rather than push them into the margin.
\sloppy
\emergencystretch=3em

\begin{document}
\maketitle
\vspace{-8mm}

\begin{keybox}{What this guide covers}
How a conceptual-design sizing loop works, in five different pictures of the
same thing: an algorithm, a flowchart, an XDSM, a design structure matrix, and
a fixed-point plot. It then compares the two sizing loops in the course
framework, \texttt{SizingLoopL1} and \texttt{SizingLoopL2}, on the same
aircraft, with numbers.

Companion material: \texttt{lecture/L01}--\texttt{L07} (Boeing 777-200LR,
instructor walkthrough) and \texttt{lab/LAB00}--\texttt{LAB01} (F-16A, your
turn). Every number quoted here is reproduced by those files.

Every diagram in this guide is also in \texttt{guide/diagrams.html}, where you
can zoom into it. Useful for the tall flowcharts.
\end{keybox}

\tableofcontents
\vspace{4mm}

% =========================================================================
\section{The sizing problem}

Sizing answers one question: \emph{what is the smallest aircraft that can fly
this mission and meet these requirements?} Three numbers come out of it.

\begin{center}
\begin{tabular}{@{}lll@{}}
\toprule
Symbol & Meaning & Units \\
\midrule
$W_{TO}$  & takeoff gross weight & lbf \\
$S_{ref}$ & wing reference area  & ft$^2$ \\
$T_{SL}$  & sea-level static thrust & lbf \\
\bottomrule
\end{tabular}
\end{center}

\noindent Everything else in conceptual design hangs off those three.

The difficulty is circularity. Start from the only statement you can write
down with confidence, which is bookkeeping:
\begin{equation}
W_{TO} = W_{empty} + W_{fuel} + W_{payload}
\label{eq:book}
\end{equation}
Payload is a requirement, so it is known. The other two are not, and worse,
both of them \emph{depend on $W_{TO}$}:

\begin{itemize}[nosep]
  \item a heavier aircraft needs a bigger wing and a bigger engine, so it has
        more structure, so $W_{empty}$ rises;
  \item a heavier aircraft makes more drag at the same lift coefficient, so it
        burns more fuel over the same mission, so $W_{fuel}$ rises.
\end{itemize}

\noindent You cannot evaluate the right-hand side of Eq.~\ref{eq:book} without
already knowing the left-hand side. That is what makes sizing a loop rather
than a calculation.

\begin{keybox}{The one sentence to remember}
Sizing is a fixed-point problem. Guess a takeoff weight, ask the discipline
models what that aircraft would actually weigh, and keep going until the
answer equals the question.
\end{keybox}

% =========================================================================
\section{The three governing relations}

The loop is built from three pieces. Each one is a separate body of theory,
and each one is a separate object in the code.

\subsection{Closure: the weight equation}

Divide Eq.~\ref{eq:book} through by $W_{TO}$ and solve for the takeoff weight
the payload requires:
\begin{equation}
\boxed{\;W_{TO}^{new} = \frac{W_{payload}}
       {1 - \dfrac{W_{empty}}{W_{TO}} - \dfrac{W_{fuel}}{W_{TO}}}\;}
\label{eq:closure}
\end{equation}

\noindent Reference: Raymer, \emph{Aircraft Design: A Conceptual Approach},
6th ed., Eq.~3.4; Martins design metabook, Algorithm~1. In the code this is
\texttt{SizingSteps.togw\_update}, and the denominator is returned alongside
the answer because it carries so much information.

\begin{warnbox}{Watch the denominator}
If the empty fraction and the fuel fraction together reach 1, the denominator
hits zero and no positive $W_{TO}$ closes the payload. The function returns
\texttt{NaN} and the sizing loops raise
\texttt{SizingLoopL1:closureInfeasible}. Physically this says: this aircraft,
at this design point, is all structure and fuel and has nothing left to carry.
\end{warnbox}

\subsection{The design point: constraint analysis}

Eq.~\ref{eq:closure} gives a weight, not a shape. The shape comes from the
performance requirements. Each requirement becomes one curve in the
thrust-to-weight against wing-loading plane. Most of them are \emph{producers}
built on the Mattingly master equation,
\begin{equation}
\frac{T}{W} = \frac{A}{W/S} + B\,\frac{W}{S} + C + D,
\qquad
\begin{aligned}
A &= \frac{q\,C_{D0}}{\alpha}, &
B &= \frac{q}{\alpha}K_1\!\left(\frac{n\beta}{q}\right)^{\!2},\\
C &= \frac{K_2\,n\,\beta}{\alpha}, &
D &= \frac{\beta}{\alpha}\,\frac{P_s}{V},
\end{aligned}
\label{eq:master}
\end{equation}
and a few are \emph{walls}, which cap $W/S$ and say nothing about thrust.
Landing field length is the usual wall. Reference: Mattingly,
\emph{Aircraft Engine Design}, 2nd ed.; Raymer ch.~5.

A feasible design sits above every producer and left of every wall. The
\emph{design point} $\left((W/S)^*, (T/W)^*\right)$ is the best corner of that
region: for a jet, the smallest engine, so the lowest $T/W$, and then the
highest $W/S$ that still reaches it. In the code,
\texttt{ConstraintAnalysis.optimal\_point\_continuous} finds it with
\texttt{fmincon}, minimising $T/W$ subject to every signed residual
$g_i = \text{required} - \text{available} \le 0$.

Given the design point, a weight fixes both sizes:
\begin{equation}
S_{ref} = \frac{W_{TO}}{(W/S)^*}, \qquad T_{SL} = (T/W)^* \, W_{TO}.
\label{eq:slaved}
\end{equation}

\begin{keybox}{The most useful output of constraint analysis}
\texttt{[WS, TW, info] = con.optimal\_point\_continuous()} returns
\texttt{info.active\_names}: the requirements whose residual is exactly zero at
the optimum. Those are the ones holding the design where it is. Everything else
is satisfied with margin and is not driving anything.

For the Boeing 777 in this package there is exactly one:
\emph{Climb 3 (2nd segment, FAR 25.121)}. For the F-16A at Level 2 there are
two: \emph{Max Mach} and \emph{Landing}.
\end{keybox}

\subsection{The fuel fraction: mission analysis}

The mission model walks the flight profile leg by leg, threading the weight
through. Cruise and loiter legs use Breguet; startup, taxi, takeoff, climb,
descent and landing use historical fractions (Roskam Part~I Table~2.1). A
reserve markup is applied to the total. In the code,
\texttt{[W\_fuel, breakdown] = miss.total\_fuel(W\_TO)}.

Note that the mission profile lives in the \textbf{requirements} file, not in
any discipline model. That separation matters: a \emph{spec} file says what the
aircraft \emph{is}; a \emph{requirements} file says what it must \emph{do}.
Sizing changes the spec until it satisfies the requirements.

% =========================================================================
\section{Algorithm form}

\begin{algorithm}[H]
\SetAlgoLined
\DontPrintSemicolon
\KwIn{$W_{TO}^{guess}$; tolerance $\varepsilon = 10^{-6}$; relaxation $w = 0.5$}
\KwOut{$W_{TO}$, $S_{ref}$, $T_{SL}$}
$\left((W/S)^*, (T/W)^*\right) \leftarrow$ \texttt{con.optimal\_point\_continuous()}
   \tcp*{ONCE, outside the loop}
$W \leftarrow W_{TO}^{guess}$\;
\For{$k = 1$ \KwTo $k_{max}$}{
  $S_{ref} \leftarrow W / (W/S)^*$ \tcp*{write into geom}
  $T_{SL}  \leftarrow (T/W)^* \cdot W$ \tcp*{write into prop}
  $W_{fuel} \leftarrow$ \texttt{miss.total\_fuel}$(W)$\;
  $W_{OEW}  \leftarrow$ \texttt{wts.get\_OEW}$(W)$\;
  $\mathrm{denom} \leftarrow 1 - W_{OEW}/W - W_{fuel}/W$\;
  \lIf{$\mathrm{denom} \le 0$}{\textbf{error}: the closure is infeasible}
  $W^{new} \leftarrow W_{payload} / \mathrm{denom}$ \tcp*{Raymer Eq.\ 3.4}
  \lIf{$|W^{new} - W| / W^{new} < \varepsilon$}{$W \leftarrow W^{new}$; \textbf{converged}; \textbf{break}}
  $W \leftarrow W + w\,(W^{new} - W)$ \tcp*{under-relax}
}
\caption{\texttt{SizingLoopL1} --- one state variable. [Martins slide 6]}
\end{algorithm}

\vspace{3mm}

\begin{algorithm}[H]
\SetAlgoLined
\DontPrintSemicolon
\KwIn{$W_{TO}^{guess}$, $T_{SL}^{guess}$; $\varepsilon = 10^{-6}$; $w_W = w_T = 0.5$}
\KwOut{$W_{TO}$, $T_{SL}$, $S_{ref}$, $S_{ht}$, $S_{vt}$}
$W \leftarrow W_{TO}^{guess}$; \; $T \leftarrow T_{SL}^{guess}$\;
\texttt{prop.T\_SL} $\leftarrow T$ \tcp*{seed BEFORE the first solve}
$\left((W/S)^*, (T/W)^*\right) \leftarrow$ \texttt{con.optimal\_point\_continuous()}\;
\For{$k = 1$ \KwTo $k_{max}$}{
  $S_{ref} \leftarrow W / (W/S)^*$ \tcp*{1. wing}
  $(S_{ht}, S_{vt}) \leftarrow$ \texttt{tail.size()} \tcp*{2. tail, before the weights}
  \texttt{prop.T\_SL} $\leftarrow T$ \tcp*{3. thrust, before the solve}
  $\left((W/S)^*, (T/W)^*\right) \leftarrow$ \texttt{optimal\_point\_continuous}$([\,(W/S)^*, (T/W)^*\,])$ \tcp*{4. EVERY pass}
  $T^{new} \leftarrow (T/W)^* \cdot W$\;
  $W_{fuel} \leftarrow$ \texttt{miss.total\_fuel}$(W)$ \quad
  $W_{OEW} \leftarrow$ \texttt{wts.get\_OEW}$(W)$ \tcp*{5.}
  $W^{new} \leftarrow W_{payload} / (1 - W_{OEW}/W - W_{fuel}/W)$ \tcp*{6.}
  \lIf{$|W^{new}-W|/W^{new} < \varepsilon$ \textbf{and} $|T^{new}-T|/T^{new} < \varepsilon$}{\textbf{converged}; \textbf{break}}
  $W \leftarrow W + w_W (W^{new} - W)$ \quad
  $T \leftarrow T + w_T (T^{new} - T)$\;
}
\caption{\texttt{SizingLoopL2} --- two state variables. [Martins slide 8]}
\end{algorithm}

\begin{warnbox}{Three details in Algorithm 2 that are easy to get wrong}
\begin{enumerate}[nosep, leftmargin=*]
  \item \textbf{Order.} Wing, then tail, then thrust, \emph{then} the design-point
        solve. Sizing the tail after the weights means the weights read last
        pass's tail. Writing the thrust after the solve means the solve reads
        last pass's drag.
  \item \textbf{$T^{new} = (T/W)^* W$.} The new thrust demand comes from the
        design point and the current weight, not from the relaxed state $T$.
  \item \textbf{The convergence test is \emph{and}, not \emph{or}.} Stopping when
        the weight settles leaves the thrust still moving.
\end{enumerate}
\end{warnbox}

\subsection{Under-relaxation}

Both loops end the pass with
\begin{equation}
x \leftarrow x + w\,(x^{new} - x),
\label{eq:relax}
\end{equation}
which is \texttt{SizingSteps.relax(x\_old, x\_new, w)}.

\begin{warnbox}{The argument order}
$w$ weights the \textbf{new} value. $w = 1$ takes the full undamped step;
$w = 0.5$, the default, goes halfway. Writing
$x \leftarrow w\,x + (1-w)\,x^{new}$ is the opposite convention and is the most
common mistake in this lab.
\end{warnbox}

% =========================================================================
\section{Flowchart: the control flow}

The algorithm as a flowchart, including the branches and the two ways each
loop can fail. Section~\ref{sec:l1l2} explains the differences; here just note
that Level 2 has more inside the box.

\begin{figure}[p]
\centering
\includegraphics[height=0.87\textheight]{flow_L1.png}
\caption{\texttt{SizingLoopL1}. One state variable. The design point is solved
before the loop and never revisited.}
\end{figure}

\begin{figure}[p]
\centering
\includegraphics[height=0.87\textheight]{flow_L2.png}
\caption{\texttt{SizingLoopL2}. Two state variables, a tail-sizing step, and
the design point re-solved inside the loop. Yellow marks what Level 1 does not
have.}
\end{figure}

\clearpage

% =========================================================================
\section{XDSM: the data flow}

A flowchart shows what happens next. It does not show what data goes where, and
in a coupled problem that is the thing you need. The XDSM does.

\begin{keybox}{Reading an XDSM}
\emph{Lambe and Martins, ``Extensions to the Design Structure Matrix\ldots'',
Structural and Multidisciplinary Optimization 46 (2012) 273--284.}
\begin{itemize}[nosep, leftmargin=*]
  \item Cells on the diagonal are \textbf{components}: something that computes.
  \item Cells off the diagonal are \textbf{data}. Row $i$, column $j$ carries
        data produced by component $i$ and consumed by component $j$.
  \item \textbf{Above} the diagonal is feed-forward. \textbf{Below} the diagonal
        is feedback, drawn here in red.
  \item The grey highway is the data path: out of a component along its row,
        then along the consumer's column.
  \item The number at the corner of a component is its place in the process.
\end{itemize}
\end{keybox}

\begin{landscape}
\begin{figure}[p]
\centering
\includegraphics[width=\linewidth, height=0.90\textheight, keepaspectratio]{xdsm_L1.png}
\caption{XDSM of \texttt{SizingLoopL1}. One red edge: the weight closure.}
\end{figure}

\begin{figure}[p]
\centering
\includegraphics[width=\linewidth, height=0.90\textheight, keepaspectratio]{xdsm_L2.png}
\caption{XDSM of \texttt{SizingLoopL2}. One solver, three couplings: two
\emph{states} (solid red, $W_{TO}$ and $T_{SL}$, returning to the loop) and one
\emph{coupling} (dashed amber, $W/S$, returning to geometry --- closed by the
same loop, not a second solver).}
\end{figure}
\end{landscape}

% =========================================================================
\section{Design structure matrix: where the feedback is}

The XDSM carries the data labels. Strip them away and you get a plain
$n \times n$ design structure matrix, which answers one question very fast:
\emph{how many feedback couplings are there, and where?}

\begin{figure}[htbp]
\centering
\includegraphics[width=\linewidth]{n2_L1_L2.png}
\caption{Design structure matrices for the two loops. Below-diagonal marks are
iteration, all closed by one solver: red returns to the loop (a state), amber
returns to a component (a coupling). Level 1 has one red; Level 2 has two red
plus one amber.}
\end{figure}

\begin{keybox}{Why this diagram is worth drawing}
A process with no feedback is not a loop, it is a checklist: run the components
in order and you are done. Every mark below the diagonal is a place where a
later component sends something back to an earlier one, which one pass cannot
satisfy. Count the marks below the diagonal and you have counted the reasons
the process must iterate.
\end{keybox}

\clearpage

% =========================================================================
\section{One pass, as a sequence}

\begin{figure}[htbp]
\centering
\includegraphics[width=0.95\linewidth]{sequence_L2.png}
\caption{One pass of \texttt{SizingLoopL2}, in the order the calls actually
happen. Note how much of the coupling travels through shared handles rather
than through arguments.}
\end{figure}

\begin{warnbox}{Handle semantics}
Every discipline object is a MATLAB \texttt{handle}. The sizing loop does not
return a new aircraft; it \textbf{writes into the objects you gave it}. Derived
geometry properties are \texttt{Dependent} and recompute on every read, so a
change to \texttt{geom.S\_ref} is visible to the aerodynamics model
immediately, with no argument passing anything.

The practical consequence: \textbf{build a fresh stack before every run}. A
second run on a used stack starts from the first run's leftovers.
\end{warnbox}

% =========================================================================
\section{The fixed point}

Everything above is one equation, $W = f(W)$, where $f$ is a single closure
pass. Plot $f$ against the 45-degree line and the whole method becomes visible.

\begin{figure}[htbp]
\centering
\includegraphics[width=0.72\linewidth]{cobweb_b777.png}
\caption{The Boeing 777 closure map. Where $f$ crosses the 45-degree line,
input equals output and the aircraft is sized. The staircase is the iteration.
Produced by \texttt{lecture/L04}.}
\end{figure}

The slope of $f$ at the crossing carries two separate meanings.

\paragraph{It decides whether the iteration converges.} An error $e$ becomes
$f'e$ after one plain pass. If $|f'| < 1$ the error shrinks; if $|f'| > 1$ it
grows and plain substitution diverges. Under-relaxation changes the effective
slope to
\begin{equation}
f'_{\text{eff}} = 1 + w\,(f' - 1),
\end{equation}
which is how a smaller $w$ rescues a steep map.

\paragraph{It is the growth factor.} $f'$ is also how many pounds of takeoff
weight one extra pound of payload buys. See Section~\ref{sec:growth}.

\begin{keybox}{A worked case: the 777 diverges without relaxation}
\texttt{lecture/L04} runs plain substitution on the Boeing 777. It does not
converge. The first eight passes go
\[
700{,}000 \rightarrow 822{,}000 \rightarrow 630{,}000 \rightarrow 1{,}066{,}000
\rightarrow 480{,}000 \rightarrow 10{,}349{,}000 \rightarrow 272{,}000
\rightarrow \texttt{NaN}.
\]
The slope estimated from the first two passes is
$(629{,}500 - 822{,}400)/(822{,}400 - 700{,}000) \approx -1.58$: steeper than
1 in magnitude, and negative, so the error grows and flips sign every pass. By
pass 7 the denominator has gone negative and there is no such aircraft.

At $w = 0.5$ the effective slope becomes $1 + 0.5(-1.58 - 1) = -0.29$, and the
same problem converges in 14 passes.
\end{keybox}

\clearpage

% =========================================================================
\section{Level 1 against Level 2}
\label{sec:l1l2}

\subsection{Two different things are called ``Level 2''}

\begin{figure}[htbp]
\centering
\includegraphics[width=\linewidth]{fidelity_ladder.png}
\caption{The two axes. Nothing forces them to agree.}
\end{figure}

\textbf{Loop fidelity} is how many unknowns the sizing loop solves for and what
it re-solves on every pass: \texttt{SizingLoopL1} against
\texttt{SizingLoopL2}. \textbf{Discipline fidelity} is how one model computes
its own answer: a statistical regression, a component build-up, or a detailed
physical geometry.

The Boeing 777 example in this package runs a \emph{Level-2 loop} over
\emph{Level-2 geometry and weights} but \emph{Level-1 aerodynamics and
propulsion}. That mix is deliberate, because for a transport it is the empty
weight that drives the answer.

There is no \texttt{SizingLoopL3}. Sizing has no equations of its own that
change with fidelity, only a count of unknowns, so the same
\texttt{SizingLoopL2} class drives the Level-3 F-16 rung as well.

\subsection{What changes in the loop}

\begin{center}
\footnotesize
\begin{tabular}{@{}p{42mm}p{48mm}p{62mm}@{}}
\toprule
& \textbf{\texttt{SizingLoopL1}} & \textbf{\texttt{SizingLoopL2}} \\
\midrule
Injected objects & 6 & \textbf{7} --- adds \texttt{tail} \\
State variables & $W_{TO}$ & \textbf{$W_{TO}$ and $T_{SL}$} \\
Design point $\left((W/S)^*,(T/W)^*\right)$
  & \textbf{frozen}: solved once, before the loop, and not even re-solved after it
  & \textbf{freed}: re-solved every pass, warm-started from the last one, plus once more at the end \\
$T_{SL}$ inside the loop
  & slaved, $(T/W)^* W$, write-only, nothing reads it back
  & an independent under-relaxed state; read by geometry (nacelle $\rightarrow$ $C_{D0}$), by weights (engine weight), and by the mission \\
Tail & never sized & \texttt{tail.size()} every pass $\rightarrow$ \texttt{geom.S\_ht}, \texttt{geom.S\_vt} \\
Mission breakdown & discarded & consumed $\rightarrow$ \texttt{wts.W\_landing} (Level 3 only) \\
Convergence test & $|\Delta W|/W^{new} < \varepsilon$ & the same \textbf{and} the same test on $T_{SL}$ \\
\texttt{history} fields & 6 & 13 \\
Feedback couplings & 1 & \textbf{3}: two states $+$ one coupling \\
Citation & Martins slide 6 & Martins slide 8 \\
\bottomrule
\end{tabular}
\end{center}

\subsection{Why Level 1 is allowed to freeze the design point}

Because at Level 1 the aerodynamics object takes \textbf{no geometry}. It reads
aspect ratio and sweep as plain numbers from the spec file. Nothing the loop
changes can move a drag polar, so nothing can move a constraint curve, so the
design point genuinely cannot change.

At Level 2 the aerodynamics object \emph{holds} the geometry object, and
$C_{D0} = C_{fe} S_{wet} / S_{ref}$. Both of the things the loop writes ---
the wing area and the thrust, through the nacelle --- move the wetted area.
Measured on the 777 in \texttt{lecture/L06}: shrink the wing by a quarter and
$C_{D0}$ goes from $0.01597$ to $0.01956$, moving the design point from
$(161\ \mathrm{psf},\ 0.2657)$ to $(165\ \mathrm{psf},\ 0.2722)$.

\subsection{What it actually buys, measured}

Same aircraft, same discipline models, same starting guess. From
\texttt{lecture/L06}:

\begin{center}
\footnotesize
\begin{tabular}{@{}lrrr@{}}
\toprule
Boeing 777-200LR & \textbf{Level-1 loop} & \textbf{Level-2 loop} & Real aircraft \\
\midrule
$W_{TO}$ [lbf]     & 740{,}125 & 740{,}060 & 766{,}800 \\
$T_{SL}$ [lbf]     & 196{,}671 & 196{,}674 & 220{,}000 \\
$S_{ref}$ [ft$^2$] & 4{,}597   & 4{,}598.5 & 4{,}605 \\
$(W/S)^*$ [psf]    & 161.00    & 160.93    & 142.45 \\
$(T/W)^*$          & 0.26573   & 0.26575   & 0.28691 \\
$S_{ht}$ [ft$^2$]  & \emph{never sized} & 907.8 & --- \\
$S_{vt}$ [ft$^2$]  & \emph{never sized} & 800.7 & --- \\
passes             & 14        & 52        & --- \\
\bottomrule
\end{tabular}
\end{center}

Three things to take from that table.

\paragraph{1. The answer barely moved.} Less than a hundredth of a percent on
takeoff weight. Say this out loud in class: extra fidelity did \emph{not} change
the answer here. The reason is visible in the spec file --- the wing area
already in it was 4{,}605 ft$^2$ and the loop converged to 4{,}598 ft$^2$. The
starting aircraft was already the finished aircraft, so Level 1 froze its design
point in the right place and got away with it.

\begin{warnbox}{Sized is not the same as consumed}
It is tempting to go on and say \emph{and the empty weight sees the new tail}.
On this aircraft that is \textbf{false}. \texttt{B777GeomL2} derives its exposed
tail areas from the fixed Chapter-7 trapezoids, not from the loop-written
reference areas, so on the 777 the tail resize is \textbf{write-only}. Double
\texttt{geom.S\_ht} and \texttt{S\_exposed\_ht}, \texttt{S\_wet},
\texttt{W\_tail} and \texttt{get\_OEW} all move by exactly zero.

On the \textbf{F-16A} the same test moves every one of them, because
\texttt{F16GeomL2} builds its tail chords and span from \texttt{S\_ht}. Same
loop, same step, different coupling, because the two geometry models were
built to different plans. Section~\ref{sec:trace} shows both.
\end{warnbox}

\paragraph{2. Level 2 sized the tail; Level 1 could not.} There is no tail
object in \texttt{SizingLoopL1}, and the 777 geometry file does not carry tail
areas, so after a Level-1 run they are still undefined.

\paragraph{3. Now start from the wrong aircraft.} Put 2{,}600 ft$^2$ in the
spec instead of 4{,}605 and change nothing else:

\begin{center}
\footnotesize
\begin{tabular}{@{}lrrrr@{}}
\toprule
& L1, good start & L1, \textbf{bad start} & L2, good start & L2, \textbf{bad start} \\
\midrule
$W_{TO}$ [lbf]  & 740{,}125 & \textbf{761{,}614} & 740{,}060 & 739{,}821 \\
$(W/S)^*$ [psf] & 161.00    & \textbf{170.00}    & 160.93    & 160.75 \\
$(T/W)^*$       & 0.26573   & \textbf{0.28068}   & 0.26575   & 0.26574 \\
\bottomrule
\end{tabular}
\end{center}

\begin{keybox}{What Level 2 is really for}
Level 1 moved by 2.9\% and reported a design point that is plainly wrong: 170
psf and 0.281 against a true optimum near 161 psf and 0.266. Level 2 landed
within a few hundred pounds of where it was before.

So the value of the Level-2 loop is not a better number on a good day. It is
\textbf{not depending on the guess you started from}. In real design work you
do not know the answer in advance, which is exactly the situation the bad-start
column reproduces.

The cost is real: 52 passes against 14, each one running \texttt{fmincon}
instead of reusing a frozen answer.
\end{keybox}

\clearpage

% =========================================================================
\section{What actually moves, pass by pass}
\label{sec:trace}

\subsection{Two unknowns, three outputs}

Sizing produces three numbers, so it is natural to assume the loop iterates on
three unknowns. It does not.

\begin{center}
\begin{tabular}{@{}lll@{}}
\toprule
Quantity & Kind & How it is produced \\
\midrule
$W_{TO}$  & \textbf{STATE}  & under-relaxed each pass from the weight closure \\
$T_{SL}$  & \textbf{STATE}  & under-relaxed each pass from $(T/W)^*\,W_{TO}$ \\
$S_{ref}$ & \textbf{SLAVED} & recomputed, never iterated: $S_{ref} = W_{TO}/(W/S)^*$ \\
\bottomrule
\end{tabular}
\end{center}

All three \emph{change} on every pass, but only two are unknowns.
\texttt{SizingLoopL2} under-relaxes exactly two quantities and tests exactly two
residuals. The wing area moves for two reasons at once: the weight it carries
changes, \emph{and} the design point it is sized against changes. The second
reason is the one Level 1 does not have.

\subsection{Tracing the loop}

\texttt{SizingLoopL2} logs 13 quantities per pass, which is enough to know
whether it converged and not enough to know why. The package ships
\texttt{helpers/trace\_sizing\_L2.m}, a replica of the loop body that records
about thirty. It checks itself against the class at the end of every call and
errors if the two have drifted, so the trace is the real loop.

\begin{lstlisting}
T = trace_sizing_L2(@build_f16_stack_L2, 30000, 20000);
plot_sizing_history(T);      % twelve panels
plot_weight_history(T);      % the component build-up, pass by pass
plot_what_changed(T);        % one bar chart: what moved, what did not
\end{lstlisting}

\begin{figure}[htbp]
\centering
\includegraphics[width=\linewidth]{history_f16.png}
\caption{F-16A, Level-2 loop, 14 passes. Panel 4 is the one to look at: if the
design point were flat there, a Level-1 loop would have been good enough. Here
$(T/W)^*$ moves 24\%.}
\end{figure}

\subsection{Smaller is draggier}

Follow the chain in panels 7 to 10. The wing shrinks, so the exposed and wetted
areas shrink --- but $S_{ref}$ shrinks \emph{faster} than $S_{wet}$ does,
because the fuselage and the inlet duct do not shrink at all. The ratio
$S_{wet}/S_{ref}$ therefore \textbf{rises}, 5.09 to 6.18. Since
$C_{D0} = C_{fe} S_{wet}/S_{ref}$, parasite drag rises 21\%, $L/D_{max}$ falls
from 11.0 to 9.9, every constraint curve lifts, and the design point demands
more thrust.

That is a real conceptual-design effect and it is worth a minute of class time:
\textbf{a smaller aeroplane is a draggier aeroplane, per unit of wing.} It is
one reason fighters do not shrink indefinitely.

\begin{figure}[htbp]
\centering
\includegraphics[width=0.92\linewidth]{weights_f16.png}
\caption{The empty-weight build-up. Wing and tail track the areas the loop is
changing; gear and \emph{all else} track gross weight; the fuselage does not
move at all, because it is a fixed input at this fidelity.}
\end{figure}

\begin{keybox}{A flat band is information}
A component that does not move is one the loop has no lever on. Any error in it
passes straight through to $W_{OEW}$ without being traded against anything
else. Knowing which components are inert is as useful as knowing which are
active.
\end{keybox}

\clearpage

\subsection{One chart per aircraft}

Every traced quantity, as a percentage change from the first pass to the last,
coloured by what kind of thing it is.

\begin{figure}[htbp]
\centering
\includegraphics[width=0.82\linewidth]{changed_f16.png}
\caption{F-16A. Everything is coupled: the tail areas move, and so do the
exposed tail area and the tail weight.}
\end{figure}

\begin{figure}[htbp]
\centering
\includegraphics[width=0.82\linewidth]{changed_b777.png}
\caption{B777-200LR. \texttt{S\_ht} and \texttt{S\_vt} move $+8.8\%$, but
\texttt{S\_exp\_ht}, \texttt{S\_exp\_vt} and \texttt{W\_tail} all read
\emph{did not move}. On this aircraft the tail resize is write-only.}
\end{figure}

\begin{warnbox}{Trace before you trust}
Both aircraft converge. Both look healthy. Only one of them has the tail
coupling you would have assumed from reading the loop. The converged answer
cannot tell you this; the trace can.
\end{warnbox}

\clearpage

% =========================================================================
\section{Growth factor}
\label{sec:growth}

Add one pound of payload to a sized aircraft and the takeoff weight goes up by
more than one pound, because the loop settles at a new and heavier fixed point.
The amplification is the \textbf{growth factor}, and it is the most useful
single number conceptual design produces.

If the empty and fuel fractions held still, the growth factor would be exactly
$1/\mathrm{denom}$. They do not hold still. For the 777, $1/\mathrm{denom} =
9.4$, but re-sizing with 5{,}000 lbf more payload gives a measured growth
factor of \textbf{3.3}. The gap is real physics, not a mistake: follow the
\texttt{denom} column in the \texttt{lecture/L04} table and it \emph{rises} with
aircraft size (0.074 at 630{,}000 lbf, 0.107 at 740{,}000, 0.164 at
1{,}070{,}000). Each pound of growth buys a little less growth than the pound
before it, so the runaway damps itself.

Keep $1/\mathrm{denom}$ anyway, as a warning light. A small denominator means
structure and fuel have eaten nearly the whole aircraft, and such a design is
violently sensitive to everything.

\subsection{Which requirements are expensive}

\texttt{lecture/L07} nudges four inputs by $\pm 10\%$ and re-sizes each time.
Dividing by the input change gives an \emph{elasticity}: percent change in
$W_{TO}$ per percent change in the input.

\begin{center}
\footnotesize
\begin{tabular}{@{}lrrr@{}}
\toprule
Input (Boeing 777) & $W_{TO}$ at $-10\%$ & $W_{TO}$ at $+10\%$ & elasticity (up) \\
\midrule
cruise TSFC            & $-9.66\%$ & $+10.38\%$ & \textbf{1.04} \\
cruise range           & $-9.42\%$ & $+10.09\%$ & \textbf{1.01} \\
skin friction $C_{fe}$ & $-4.79\%$ & $+4.67\%$  & 0.47 \\
payload                & $-3.57\%$ & $+3.47\%$  & 0.35 \\
\bottomrule
\end{tabular}
\end{center}

\begin{keybox}{Reading the table}
Cruise range and cruise TSFC both come out near 1.0. A one percent longer
mission, or a one percent thirstier engine, buys a one percent heavier
aeroplane. Those are the expensive knobs, and they are the ones worth an
argument with the customer or with the engine supplier.

Payload comes out near 0.35, which does \emph{not} contradict the growth factor
of 3.3. Each \emph{pound} of payload still costs 3.3 lbf of aircraft, but
payload is only about 11\% of the takeoff weight, so ten percent of it is not
many pounds. Growth factor and elasticity answer different questions --- one per
pound, one per percent --- and design work needs both.

Note also that technology knobs (TSFC, $C_{fe}$) and requirement knobs (payload,
range) land on the same axis in the same units. That is how conceptual design
decides whether to buy a better engine or to renegotiate the mission.
\end{keybox}

% =========================================================================
\section{Where each equation lives in the code}

\begin{center}
\footnotesize
\begin{tabular}{@{}p{42mm}p{54mm}p{50mm}@{}}
\toprule
\textbf{What} & \textbf{File} & \textbf{Reference} \\
\midrule
Weight closure, under-relaxation
  & \texttt{src/sizing/SizingSteps.m}
  & Raymer 6th ed.\ Eq.~3.4; metabook Algorithm~1 \\
One-state loop
  & \texttt{src/sizing/SizingLoopL1.m}
  & Martins slide 6 \\
Two-state loop
  & \texttt{src/sizing/SizingLoopL2.m}
  & Martins slide 8 \\
Design point, \texttt{fmincon}
  & \texttt{src/constraints/ConstraintAnalysis.m}
  & Raymer ch.~5 \\
Master equation curves
  & \texttt{src/constraints/MasterEquationConstraint.m}
  & Mattingly 2nd ed. \\
Mission fuel
  & \texttt{src/core/mission/}
  & Roskam Part~I Eqs.~2.10--2.15 \\
Tail volume coefficients
  & \texttt{src/disciplines/tail\_sizing/TailL1.m}
  & Raymer 7th ed.\ Table~6.4 \\
Standard atmosphere
  & \texttt{src/core/AircraftState.m}
  & ICAO 1993, via \texttt{atmosisa} \\
\bottomrule
\end{tabular}
\end{center}

\begin{keybox}{Toolboxes}
The sizing path needs two: \textbf{Aerospace Toolbox} for \texttt{atmosisa},
which builds every flight condition, and \textbf{Optimization Toolbox} for
\texttt{fmincon}, which both loops use to find the design point.
\texttt{START\_HERE.mlx} checks for both.
\end{keybox}

% =========================================================================
\section{Running the material}

\begin{center}
\footnotesize
\begin{tabular}{@{}p{52mm}p{100mm}@{}}
\toprule
\texttt{START\_HERE.mlx} & path setup, toolbox check, and a smoke solve \\
\midrule
\multicolumn{2}{@{}l}{\textbf{\texttt{lecture/} --- Boeing 777-200LR, instructor walkthrough}} \\
\texttt{L01\_build\_the\_stack}      & the six objects, and why the construction order is forced \\
\texttt{L02\_constraint\_design\_point} & ten requirements $\rightarrow$ one design point \\
\texttt{L03\_fuel\_and\_empty\_weight}  & the two questions the closure asks \\
\texttt{L04\_one\_closure\_step\_by\_hand} & one pass, every number on screen; divergence; under-relaxation \\
\texttt{L05\_run\_the\_L1\_loop}     & the Level-1 loop to convergence \\
\texttt{L06\_run\_the\_L2\_loop}     & the Level-2 loop, and the bad-start test \\
\texttt{L07\_sensitivity\_growth\_factor} & growth factor and elasticities \\
\midrule
\multicolumn{2}{@{}l}{\textbf{\texttt{lab/} --- F-16A, your turn}} \\
\texttt{LAB00\_warmup\_L1\_loop}     & complete and runnable; read every line \\
\texttt{LAB01\_write\_the\_L2\_loop} & six \texttt{TODO}s; a checker grades your loop against \texttt{SizingLoopL2} \\
\texttt{LAB\_worksheet.pdf}          & experiments to run on the loop you wrote \\
\bottomrule
\end{tabular}
\end{center}

\begin{warnbox}{On the F-16A numbers}
The lab converges to roughly 23{,}300 lbf against a Brandt F-16A reference of
31{,}377 lbf. That gap is a statement about the textbook discipline models and
the requirement set, \textbf{not} about your loop. The lab grades one thing:
does your loop reproduce \texttt{SizingLoopL2} to within a pound? Agreement with
a real aircraft is a separate question and the framework keeps the two
deliberately apart.
\end{warnbox}

\vfill
\begin{center}
\footnotesize\itshape
Every number in this guide is reproduced by the Live Scripts in this package.
\end{center}

\end{document}
